\scanpage{8}
\mspage{6}
lieu isolé, ordre 1: \unsure{lieu-bout}.

Ordre fini:

\deleted{Point rég}\deleted{\inserted{lieu}}

Lieu régulier $x$: $\exists I\in S$ t.q. $x\in\tilde I^\circ$.
\[
L_r=\text{ens des lieux réguliers}=\bigcup_{I\in S}\tilde I^\circ
\]
\[
L\setminus L_r=L_s\qquad
\text{lieux singuliers, ou critiques, ou de branchement.}
\]

$F_4)$ Soient $I',I''\in S$, $z\in L$, tels que
$\tilde I'\cap\tilde I''=\delta I'\cap\delta I''=\{z\}$, avec $z$ \uline{régulier}.
Posons $\delta I'=\{x,z\}$, $\delta I''=\{y,z\}$. Alors
$\exists!\ I\in S$ tel que $\tilde I\supset\tilde I'\cup\tilde I''$
(on n'affirme pas égalité!) et $\delta I=\{x,y\}$

[donc en vertu de $F_1$, $I'=I_{x,z}$, $I''=I_{y,z}$]
\margin{\uline{Axiome} de \uline{recollement}}

\deleted{\unsure{Prop.}}\inserted{\uline{Cor 1}}
Soit $I\in S$, $x\in\delta I\cap L_r$, \deleted{avec \unsure{$x\in\delta J$} \unsure{et}} ($x$ \uline{régulier}).
\deleted{Alors $\exists J\in S$ tel que $\tilde J\cap\tilde I=\deleted{\unsure{\delta J\cap\delta I}}\{x\}$.}
Posons $\delta I=\{x,y\}$. Alors $\exists J\in S$, tel que
$\tilde J\supset\tilde I$, et: $\delta J=\{y',y\}$ où $y'\neq x$
i.e. $\delta J\cap\delta I=\{y\}$ (et $x\in\tilde J^\circ$)

\margin{\deleted{\unsure{extension}}\inserted{\unsure{prolongement}} d'un segment à une extrémité régulière}

(On utilise $F_1,F_2$ pour conclure que l'ordre de $x$ est $\geq2$, puis $F_3)$ pour trouver $K\in S$ avec $\tilde K\cap\tilde I=\delta K\cap\delta I=\{x\}$, puis on applique $F_4$)
\mspage{7}

\rotatebox{90}{\reflectbox{\unsure{EXTRA STRONG}}}
\rotatebox{90}{\reflectbox{\unsure{EXTRA STRONG}}}
